% Section 3, rewritten in loop0009 item 04 (review change 7; the owner's notes).
% Sources: paper2/popularity.md ("For the paper"), the numbers of
% paper2.section2_figures.numbers() (OpenAlex caches of 2026-09-29), the period
% counts of paper2.trends.binned, Chu & Stuckey (2009) section 1
% (literature/chu_stuckey_2009.pdf, pp. 2-3), and the graph benchmark counts of
% paper2/solver_fix.md ("Table 1: proved widths", repaired rules) against
% Coudert, Mazauric & Nisse (2016) as recorded in
% pathwidth_solver/bench/reference.py. Tables from tables/ (make_tables.py).
% The method detail, the timeline and the full names table are Appendix C.
\section{The names and the communities}
\label{sec:names}

Section~\ref{sec:equiv} showed that eight of the twelve rows of the
Linhares--Yanasse table are one quantity. A result proved under any of their
names therefore holds under all of them. But a reader who searches under one
name will not find a result published under another. This section asks two
questions about the literature. Which name does it use for the quantity? And,
besides that name, which community has been most active? The answers are
pathwidth, by far (Figure~\ref{fig:names}), and open stacks
(Table~\ref{tab:periods}). The second answer explains where the best
exact algorithm for the family came from, and leads to Section~\ref{sec:search}.

\paragraph{How we counted.} We searched OpenAlex \citep{Priem2022} on
29 September 2026 for each problem's name and its common variants, in titles and
abstracts, within Computer Science, Mathematics, Engineering and Decision
Sciences. A search hit need not use the name in the table's sense: ``narrowness''
matches term rewriting, and ``path-width'' matches roads and tool paths. So we
read the 2,271 hits and kept the 1,593 that are about one of the twelve problems.
The 700 hits for the eleven smaller names were judged one by one, by one
annotator with model assistance; the 1,571 pathwidth hits were filtered by a
stated rule, checked on samples. A mechanical phrase rule agrees with the
judgements on 78\% of the 700 and leads to the same conclusions. The counts are
lower bounds, since a work can study a problem without naming it in its
abstract. Appendix~\ref{app:names} gives the method in full, the robustness
check and the counts per name.

\begin{figure}[H]
\centering
\includegraphics[width=\linewidth]{sec2_fig1_name_usage.pdf}
\caption{Pathwidth has more relevant works than the other eleven names
together, three times over, and three names have none. Works using each
problem's name, in title or abstract, in the sense of the Linhares--Yanasse
table; all years to September 2026. Colour gives the discipline that the table
assigns the problem.}
\label{fig:names}
\end{figure}

\subsection{Pathwidth has absorbed the family}

Figure~\ref{fig:names} gives the counts. Pathwidth has 1,213 relevant works,
21 times as many as MOSP and more than three times as many as the other eleven
names together (380). Three names have no relevant work at all: narrowness,
split bandwidth and edge separation. Two more are rare: interval thickness has
6 works and one-dimensional logic 11. Pathwidth is also the only name still
growing: its 311 works in 2020--2024 are more than the whole history of any
other name (Appendix~\ref{app:names}). So anything one would want to know about
algorithms for this class of problems was most likely published under
``pathwidth''.

\subsection{Besides pathwidth, open stacks}

\begin{table}[t]
\centering\small
\caption{Besides pathwidth, the most active name over 2005--2024 is MOSP,
although vertex separation led in 2015--2019. Relevant works per five-year
period, for the three names with the most works over 2005--2024; the next are
edge search (28) and node search (20).}
\label{tab:periods}
% Generated by python -m paper2.latex.make_tables. Source: paper2.section2_figures.numbers()["periods_2005_24"] (paper2.trends.binned on the OpenAlex trend cache).
% Do not edit by hand.
\begin{tabular}{@{}lrrrrr@{}}
\toprule
Name & 2005--09 & 2010--14 & 2015--19 & 2020--24 & 2005--24 \\
\midrule
Pathwidth & 114 & 179 & 267 & 311 & 871 \\
MOSP & 17 & 17 & 5 & 12 & 51 \\
Vertex separation & 7 & 10 & 15 & 0 & 32 \\
\bottomrule
\end{tabular}
% check: next names over 2005-24: edge search game 28, node search game 20, gate matrix layout 9

\end{table}

The VLSI names belong to the 1980s: since 2005 gate matrix layout has 9
relevant works and PLA folding 6. Table~\ref{tab:periods} shows the
three names with the most works over 2005--2024. MOSP is second, with 51,
ahead of vertex separation (32) and edge search (28). Vertex separation led
only in 2015--2019. The open-stacks community, in operations research, is thus
the one most active besides pathwidth.

\begin{table}[t]
\centering\small
\caption{The open-stacks papers are read apart from graph theory. Works citing
the papers that the Linhares--Yanasse table cites, by the discipline of the
papers they cite. Of 135 works citing an operations-research paper, 6 also cite
a graph-theory paper (4\%); of 285 citing a VLSI paper, 65 do (23\%). The 2002
paper itself counts as operations research. OpenAlex, 29 September 2026; 844
citing works in all.}
\label{tab:communities}
% Generated by python -m paper2.latex.make_tables. Source: paper2.section2_figures.numbers()["discipline_sets"] (OpenAlex citation cache, 2026-09-29).
% Do not edit by hand.
\begin{tabular}{@{}lrrrr@{}}
\toprule
& & \multicolumn{3}{c}{\dots\ of which also cite} \\
\cmidrule(l){3-5}
Works citing & All & operations research & VLSI & graph theory \\
\midrule
operations research & 135 & -- & 14 & 6 \\
VLSI & 285 & 14 & -- & 65 \\
graph theory & 507 & 6 & 65 & -- \\
\bottomrule
\end{tabular}
% check: by set of disciplines cited: GT 438, VLSI 208, OR 117, GT+VLSI 63, OR+VLSI 12, GT+OR 4, GT+OR+VLSI 2; citing works 844

\end{table}

It has also worked largely apart from graph theory. Table~\ref{tab:communities}
counts the 844 works that cite the papers of the table, by the disciplines of
the papers they cite. Only six works cite both an open-stacks paper and a
graph-theory paper. Two are by the table's authors, one is open-stacks work on
graph properties, and three are pathwidth or vertex-separation work from
computer science. The VLSI papers are read with graph theory about five times as
often.

\subsection{From an open-stacks challenge to a pathwidth solver}

It is not a coincidence that the best exact algorithm for MOSP came from the
open-stacks community and is also, on the published evidence, among the
strongest exact algorithms for pathwidth. The account is
in \citet[\S1]{ChuStuckey2009}. MOSP was the subject of the first Constraint
Modelling Challenge \citep{SmithGent2005}, posed in May 2005, which drew 13
entries. The winning entry, by García de la Banda and Stuckey, combined dynamic
programming with branch and bound and was an order of magnitude faster than
any other \citep{GarciaDeLaBandaStuckey2007}. \citet{ChuStuckey2009} built on
it. They branched on which customer's stack closes next rather than on which
pattern is made next, recorded the states already refuted, added dominance
rules, and became
five to six orders of magnitude faster on hard instances. Their solver closed
every open instance of the challenge within ten seconds. MOSP's busiest decade
in Table~\ref{tab:periods}, 2005--2014, is the decade of this work.

The customer search reads an instance only through its MOSP graph
\citep[p.~4]{ChuStuckey2009}, so it decides pathwidth
(Section~\ref{sec:search}). Ported to graphs, it proves more instances than
reported on the benchmarks where published exact pathwidth results exist, those
of \citet{CoudertMazauricNisse2016}. It proves the pathwidth of 11,194 of the
11,534 Rome graphs (97.1\%), against their 95.6\%, both within 600 seconds per
graph. It proves all 50 VSPLIB trees, against their 30, and 39 of the 73
Harwell--Boeing graphs, against their 26. These are comparisons with published
numbers obtained on other hardware, not a head-to-head run.

Two consequences shape the rest of the paper. First, a result on any of the
twelve problems should be stated, and searched for, as a result on pathwidth.
Section~\ref{sec:search} does this for the customer search. It finds that the
search's first dominance rule is, in correct form, a theorem already known in
the pathwidth literature, and that two of its published theorems are false.
Second, a benchmark for any problem in the exact core is a pathwidth benchmark,
which is how Section~\ref{sec:dataset} builds its dataset.
